\documentclass[12pt,a4paper]{article}

\usepackage{cmap}
\usepackage[T2A]{fontenc}
\usepackage[utf8]{inputenc}
\usepackage[russian]{babel}
\usepackage{amsmath,amssymb}
\usepackage[a4paper,margin=2cm]{geometry}

\setlength{\parindent}{0pt}
\setlength{\parskip}{0.45em}
\sloppy

\begin{document}

\begin{center}
{\Large\bfseries Билет 12: вопросы для проверки знаний}
\end{center}

\noindent\fbox{%
  \begin{minipage}{\dimexpr\linewidth-2\fboxsep-2\fboxrule\relax}
  \normalfont
Степенные ряды с комплексными членами. Круг и радиус сходимости. Характер
сходимости степенного ряда в круге сходимости. Формула Коши--Адамара.
Сохранение радиуса сходимости степенного ряда при почленном дифференцировании
и интегрировании ряда. Первая теорема Абеля.
  \end{minipage}%
}

\section*{Вопросы на понимание}

\begin{enumerate}
\item Зачем заменой $z=w-w_0$ сводят ряд с центром $w_0$ к ряду $\sum c_kz^k$?
\item Что при такой замене происходит с кругом сходимости исходного ряда?
\item Почему в формуле Коши--Адамара используется $\limsup$, а не обычный предел?
\item Как интерпретировать случаи $\limsup\sqrt[k]{|c_k|}=0$ и $\limsup\sqrt[k]{|c_k|}=+\infty$?
\item Почему сравнение $|z|/R$ с единицей решает вопрос о сходимости внутри и вне круга?
\item Почему внутри круга сходимости получается именно абсолютная сходимость?
\item Почему при $|z|>R$ достаточно проверить, что члены ряда не стремятся к нулю?
\item Может ли степенной ряд с одним и тем же радиусом сходиться в одной точке границы и расходиться в другой? Приведите типичный пример.
\item Почему радиус сходимости сам по себе ничего не гарантирует на границе $|z|=R$?
\item Чем поточечная сходимость при $|z|<R$ отличается от равномерной сходимости на замкнутом круге $|z|\le r<R$?
\item Почему в утверждении о равномерной сходимости оставляют зазор $r<R$?
\item Почему множители $k$ и $1/(k+1)$ не меняют радиус сходимости при почленном дифференцировании и интегрировании?
\end{enumerate}

\section*{Источники для сверки}

\texttt{target/answer\_question\_12.tex};
\texttt{IvGE\_1.pdf}, стр. 303--309.

\end{document}
